% Generated draft: requires mathematical and source review.
\documentclass[11pt]{article}
\usepackage[margin=25mm]{geometry}
\usepackage{fontspec}
\setmainfont{Latin Modern Roman}
\usepackage{amsmath,amssymb,mathtools,unicode-math}
\setmathfont{Latin Modern Math}
\usepackage{enumitem,needspace,xurl,hyperref}
\setlength{\parindent}{0pt}
\setlength{\parskip}{5pt}
\setlength{\emergencystretch}{3em}
\newcommand{\N}{\mathbb{N}}
\newcommand{\Z}{\mathbb{Z}}
\newcommand{\Q}{\mathbb{Q}}
\newcommand{\R}{\mathbb{R}}
\newcommand{\C}{\mathbb{C}}
\newcommand{\F}{\mathbb{F}}
\newcommand{\A}{\mathbb{A}}
\newcommand{\PP}{\mathbb{P}}
\newcommand{\E}{\mathbb{E}}
\newcommand{\eps}{\varepsilon}
\newcommand{\dd}{\,\mathrm{d}}
\begin{document}
\section*{UnsolvedMath ID: 11000289}
\textbf{Source problem code:} catalogue entry
\subsection{Definitions and mathematical statement}
Let \(P\) be a surface group or a finitely generated free group. For an automorphism \(\phi:P\to P\), a characteristic finite-index subgroup \(H\le P\) is preserved, and \(\phi\) induces a homological action
\[
 \phi_*:H_1(H;R)\longrightarrow H_1(H;R)
\]
for a chosen coefficient system. Question 4.1 asks whether growth can be detected by such an action. The first sentence says ``homomorphism,'' but the examples and results concern automorphisms: characteristic subgroups need not be preserved by arbitrary endomorphisms. That scope must be recovered rather than silently changed.

The source supplies no growth convention (word versus conjugacy-class length, polynomial degree, exponential rate or another invariant), coefficient ring, exact detection criterion, or whether \(H\) may depend on \(\phi\). Equality of growth rates, detection of exponential versus polynomial growth and detection of polynomial degree would be distinct targets. The source emphasizes that the original group's homology is insufficient, since exponentially growing automorphisms can act trivially there.

It reports Piggott [74] answering the question for polynomial-growth free-group automorphisms and says linear-growth surface automorphisms are easily handled, while calling the exponential cases open in both settings. These are source-era assessments. The following context describes then-ongoing Handel--Mosher [43] work: a pseudo-Anosov surface automorphism and its inverse have equal exponential rate, whereas fully irreducible free-group automorphisms need not. Here iwip conventionally means no positive power preserves the conjugacy class of a proper nontrivial free factor.

The source then compares mapping tori. For \(G=\Z^n\rtimes_\phi\Z\), with \(\phi\) an automorphism, it reports equivalence of automaticity, being a CAT\((0)\) group, and satisfying a quadratic isoperimetric inequality, reference [17]. The usual latter condition is a quadratic upper bound on the Dehn function; CAT\((0)\) group means an appropriate proper cocompact isometric action on a CAT\((0)\) space, with exact conventions requiring the missing footnote 3.

For surface-group mapping tori it reports automaticity and quadratic isoperimetric inequality for all cases in its intended range, and a known characterization of the CAT\((0)\) cases. The genus/domain hypotheses are absent and cannot be extended to every surface indiscriminately. For free-by-cyclic groups it cites Brady--Bridson--Reeves [12] for nonautomatic examples, Gersten [36] for non-CAT\((0)\) examples, and [18] for quadratic isoperimetric inequality in all the intended free-group-automorphism cases. These contextual comparisons do not define the missing homological-detection criterion or settle it.
\subsection{Short English statement}
Can growth of the specified surface/free-group maps be read from homology of a characteristic finite-index subgroup? The map type and meaning of detection are missing. The source's polynomial/linear cases, inverse-rate contrast and mapping-torus comparisons are retained as qualified context.
\subsection{Sources}
\begin{itemize}
\item Local UnsolvedMath registry, permanent ID \(11000289\), exact source code \texttt{catalogue entry}, Question 4.1, PDF page 334. No individual proposer or proposed year is supplied; the title is truncated.
\item Farb, edited volume \emph{Problems on Mapping Class Groups and Related Topics}: \url{https://www.math.uchicago.edu/~farb/papers/mcgbook.pdf}. The volume has 329 explicitly typeset numbered items across chapters, repeated local numbering distinguished by PDF page. The extract includes page 328 furniture naming M. Bridson and K. Vogtmann.
\item Internal references: Piggott [74]; Handel--Mosher [43]; abelian mapping-torus equivalences [17]; Brady--Bridson--Reeves [12]; Gersten [36]; general free-by-cyclic quadratic bound [18]. Complete entries and footnote 3 are absent.
\item Extraction pdf-font-structure, accessed 2026-07-29. Source evidence NEEDS\_REVIEW presents an open problem and calls for a current primary-literature audit. Registry status open and classification OPEN-TRIAGE; the research summary explicitly cannot verify status. The exponential-open and ongoing-work descriptions remain historical source claims.
\item Dataset attribution: Ulam AI / UnsolvedMath Contributors, CC BY 4.0. Source rights NEEDS\_REVIEW: publicly accessible author/publisher copy, redistribution terms requiring review.
\end{itemize}
\end{document}
